\section{The retained side}\label{sec:retained-side}

Keep every original source of Paper~II, including the complete sum $S_2$.
Fix $0<\sigma\le10^{-10}$ satisfying Paper~II, (1.27). Put
\begin{equation}\label{eq:retB-parameters}
 L=\log N,\qquad \tau=\sigma^9,\qquad \rho=1+\tau,
 \qquad h=h_G=\frac9{5000},\qquad \eta_L=\frac1{10000}.
\end{equation}
The choice of $\sigma$ precedes the limit.
Use $W$ from \eqref{eq:not-normalization}.
Write $R_{\mathrm{fix}}$ for \eqref{eq:source-fixed-sum}.

The comparison required in this section is
\begin{equation}\label{eq:retB-goal}
\begin{aligned}
 \frac{R_{\mathrm{fix}}+S_2}{W}
 &\ge D_{\mathrm{II}}+g_6+\overline J_6+2\overline H
       +\frac{5719}{10^6}+M_{5,\mathrm{internal}}-E'(N;\sigma),\\
 E'(N;\sigma)
 &=E_{\mathrm{van}}(N;\sigma)
   +\sum_{j\in\mathcal H}\bigl(|\Lambda^0_{j,N}|+|\Lambda_{j,N}|\bigr)
   +|\Lambda_{G,N}|,\qquad
 \mathcal H=\{4,5,7,10\}.
\end{aligned}
\end{equation}
In $\mathcal H$, indices $5$ and $10$ denote the lower $S_5$
and retained $S_{10}$ parts.
Every displayed error is nonnegative.
The notation agrees with \eqref{eq:intro-retained-error}:
\begin{equation}\label{eq:retB-Lambda-alias}
 \Lambda_{j,N}=\Lambda_N(\Delta_j),\qquad
 \Lambda_{G,N}=\Lambda_N(\Delta_{G1}).
\end{equation}
The required limit is $E'(N;\sigma)\to0$.
Proposition~\ref{prop:retB-comparison} proves \eqref{eq:retB-goal}.
Its inputs are listed in Remark~\ref{rem:retB-inputs}. The six-endpoint
lemma is optional: without it, the
same proof gives \eqref{eq:retB-goal} with $g_6$ replaced by $0$.
We display the unchanged comparisons first.
Then we specify the five incremental and four saved-array conversion terms.

\subsection{The original signed chains}

Let $\mathcal I=\mathcal I_{\mathrm{row}}\sqcup\{\mathrm{T5}\}$ index the
signed chains of $R_{\mathrm{fix}}+S_2$. Here $\mathcal I_{\mathrm{row}}$ is the
set of Paper~II's nineteen row chains. The chain $\mathrm{T5}$ is Paper~II's outer
chain for the upper part of $S_5$
\cite[(2.28), (5.10), (6.1) and \S6.1, condition~1]{PaperII}.
For $j\in\mathcal I$, let $Q_{j,k}$ include its normalized coefficient.
Paper~II uses ten stages
\cite[(6.4)--(6.4d)]{PaperII}:
\begin{equation}\label{eq:retB-telescope}
 Q_{j,0}=Q_{j,10}
       +\sum_{k=1}^{10}(Q_{j,k-1}-Q_{j,k}).
\end{equation}
The initial expression is the actual arithmetic source.
The last expression is a directed rational endpoint.
For $j=\mathrm{T5}$ the stages are those of Lemma~\ref{lem:retB-S5top}:
$Q_{\mathrm{T5},k}$ is defined there for $k\le3$, and
$Q_{\mathrm{T5},k}=Q_{\mathrm{T5},3}$ for $3\le k\le10$.
For the three joint groups, use
\begin{equation}\label{eq:retB-joint}
\begin{aligned}
 \mathcal O&\in\{S_1,S_2\text{-}T_8/4,G_1\},\qquad
 Q_{\mathcal O,9}=\sum_{j\in\mathcal O}Q_{j,9},\\
 \sum_{j\in\mathcal O}Q_{j,0}
 &=E_{\mathcal O}+(Q_{\mathcal O,9}-E_{\mathcal O})
    +\sum_{j\in\mathcal O}\sum_{k=1}^{9}
             (Q_{j,k-1}-Q_{j,k}),\\
 Q_{\mathcal O,9}-E_{\mathcal O}&\ge0.
\end{aligned}
\end{equation}
The joint signs are kept before rounding.
The $G_1$ weights are $(2,1,1)$.

\paragraph{Signed sources and parent assignment.}
Here ``parent'' records the assignment of an arithmetic occurrence and its
source-restoration share. It does not assert that each Buchstab child
is numerically bounded by the unsplit parent. Let $T_{j,r}$ be the
nonnegative plain-source terms of Paper~II, (2.19) and (3.2). Let
$\Delta_j^{\rm src}$ be its signed source correction. Then
\[
 3S_1=3(T_{1,0}-T_{1,1}+T_{1,2}-T_{1,3})
                  +3\Delta_1^{\rm src},\qquad
 S_2=T_{2,0}-T_{2,1}+T_{2,2}-T_{2,3}+\Delta_2^{\rm src}.
\]
For every third-depth S1 occurrence, $E_3\le p_3$ gives
$-3S_{\rm pl}(b,p_3)\ge-3S_{\rm pl}(b,E_3)$. For S2 we make this
replacement only on Paper~II's selected whole-$K$ blocks $\mathcal F$,
including $PF$ and $\mathcal F\cap R1$. There
$-S_{\rm pl}(b,p_3)\ge-S_{\rm pl}(b,E_3)$. On the strict complement
$\mathcal B$ keep the actual $p_3$ and Paper~II's strict-BV comparison.
The last child of S2 keeps its negative sign on all three pieces
$PF$, $\mathcal F\cap R1$, and $\mathcal B$. The strict BV piece stays in S2.

For the fixed-cutoff rows use their literal divisor domains. Their
signed sources are $-2S_3$, $-S_4$, $-S_{5,\rm low}$, $-S_7$,
$+S_8$, and $+S_9$, respectively. On each source cell, an upper sieve
bound $S\le U$ gives $cS\ge cU$ when $c<0$. A lower sieve bound
$L\le S$ gives $cS\ge cL$ when $c>0$. These comparisons keep the source
corrections and the row's allocated errors.
Paper~II, (2.28), gives the exact assignment identity
\[
 -S_5=-S_{5,\rm low}-S_{5,\rm top},\quad
 S_{5,\rm low}=\sum_{1/4\le\log_Np<t_*}S(p,N^\alpha),\quad
 S_{5,\rm top}=\sum_{t_*\le\log_Np<\tau_S}S(p,N^\alpha).
\]
The endpoint $t_*$ belongs to the upper/BV part. The split
precedes the structured sieve or Heath--Brown expansion and keeps
parent labels. S5-low and T5 both have parent S5. S5-low gets no
restoration address of its own.
For S10 the source partition
$S_{10}=S_{10,\rm Ret}+S_{10,\rm Drop}$ is exact, with
$S_{10,\rm Ret},S_{10,\rm Drop}\ge0$. Hence $+S_{10}\ge+S_{10,\rm Ret}$
before its lower sieve bound.

The joint source keeps the representation weights $(2,1,1)$.
Paper~II, (4.14), gives, with its stated exceptional-source and
small-prime allowance $E_G\ge0$,
\[
 G_1\le2S(\mathcal B_{11},N^{1/2})+
          S(\mathcal B_{12},N^{1/2})+
          S(\mathcal B_{13},N^{1/2})+E_G.
\]
Make its cut-measure replacement with the original signed correction
before applying (4.18). The four signed descendants of $-G_1$ then
have coefficients $(-1,+1,-1,+1)$ on the single weighted source.
Their last positive term is bounded from below by its retained cells.
The change from actual $a_3$ to label base $E$ keeps the signed peel
correction of Paper~II, (4.45)--(4.57). The four rows share one
joint source parent and one restoration allocation.
Thus S1, S2, S4, S5, S7, S10 and G1 are the seven parents of the
sixteen listed retained rows. S3, S8 and S9 keep their original
single-comparison parents, and T5 remains the upper part of S5.
These are sign-correct source reductions for the complete parent
chains. They do not discharge the distinct S6, bad-source or
switched-source comparisons in Section~\ref{sec:source-decomposition}.

The nineteen original row chains of $\mathcal I_{\mathrm{row}}$
are the sixteen rows of Appendix~\ref{subsec:app-retB-rows} and
Paper~II's original $S_3,S_8,S_9$ comparisons. The strict BV piece of $S_2$ lies inside
the joint $S_2$ group. The source $-S_5$ of $R_{\mathrm{fix}}$ is split at
$N^{t_*}$, $t_*=7109/28000$, before any sieve step. Its lower part is the row
$S_5$-low of $\mathcal I_{\mathrm{row}}$. Its upper part is the chain
$\mathrm{T5}$. The chains of $S_6$ and $S_{\mathrm{sw}}$ are excluded.
Thus the endpoint sum over $\mathcal I$ is
\begin{equation}\label{eq:retB-base-endpoints}
 \sum_{j\in\mathcal I} Q_{j,10}
 =B_{\mathrm{print}}-\delta_{G1}+\overline J_6+2\overline H,
 \qquad \delta_{G1}=4.98\cdot10^{-54}.
\end{equation}
This follows from Paper~II's printed thirteen-term sum
\cite[\S6.4, immediately before (6.6)]{PaperII}.
Its joint endpoint is $Q_{G1,10}=L_{G1}-\delta_{G1}$;
the printed sum uses $L_{G1}$. The term $-\overline J_5$ of
\eqref{eq:retB-Bprint} is $Q_{\mathrm{T5},10}$, by
Lemma~\ref{lem:retB-S5top}. Apart from $-\overline J_6-2\overline H$,
every other term of \eqref{eq:retB-Bprint} is an endpoint. It ends a
chain of $\mathcal I_{\mathrm{row}}$ or a joint group.
This is a finite coefficient identity.
The printed sum is
\begin{equation}\label{eq:retB-Bprint}
\begin{aligned}
 B_{\mathrm{print}}={}&\underline D_1+\underline d_2-\frac{U_{3,A}}2
 -\frac{U_4+U_{5,\mathrm{low}}+U_7}{4}\\
 &+\frac{L_{8,A}+L_{9,A}+L_{10}}4+L_{G1}
 -\overline J_5-\overline J_6-2\overline H\\
 >&1.83403742109641331345957914542194044932304169727869681073064.
\end{aligned}
\end{equation}
The letters denote the rational endpoints printed in Paper~II,
Appendix~A.2--A.6 and (A.7), (A.10), (A.13)--(A.14).
The constant $\overline J_6$ keeps its rational definition:
\begin{equation}\label{eq:retB-J6-endpoint}
 \overline J_6=1.3148581959157321533513247276782285310758770468886237619585.
\end{equation}

The upper part of $S_5$ has the following outer chain. Put
$t_*=7109/28000$, $\tau_S=57/224$, and let
\begin{equation}\label{eq:retB-J5-endpoint}
 \overline J_5=\frac14\bigl(c_5+8.84\cdot10^{-57}\bigr),\qquad
 c_5=0.03658166423843258896399704402764682632822246022922592990,
\end{equation}
which is Paper~II's directed upper endpoint for
$J_5(0)=\int_{t_*}^{\tau_S}\frac{dt}{t(1/2-t)}$
\cite[(5.16), (5.20), (A.14)--(A.14a)]{PaperII}.

\begin{lemma}[The outer chain of the upper part of $S_5$]\label{lem:retB-S5top}
Define
\[
\begin{aligned}
 Q_{\mathrm{T5},0}&=-\frac1W\sum_{t_*\le\log_Np_1<\tau_S}S(p_1,N^\alpha),&
 Q_{\mathrm{T5},1}&=-\frac{T_5}W,\\
 Q_{\mathrm{T5},2}&=-J_5(0),&
 Q_{\mathrm{T5},3}&=-\overline J_5,
\end{aligned}
\]
where the first sum is the second term of Paper~II, (2.28), and $T_5$ is
Paper~II, (5.10). Then
\[
 Q_{\mathrm{T5},0}-Q_{\mathrm{T5},1}\ge-E_{\mathrm{T5},1}(N),\qquad
 Q_{\mathrm{T5},1}-Q_{\mathrm{T5},2}\ge-E_{\mathrm{T5},2}(N),\qquad
 Q_{\mathrm{T5},2}-Q_{\mathrm{T5},3}\ge0,
\]
with $E_{\mathrm{T5},1},E_{\mathrm{T5},2}\ge0$ tending to $0$ at fixed
$\sigma$. Both are entries of $E_{\mathrm{van}}$. In particular
$Q_{\mathrm{T5},0}\ge-\overline J_5-E_{\mathrm{T5},1}(N)-E_{\mathrm{T5},2}(N)$.
\end{lemma}

\begin{proof}
Paper~II splits $-S_5$ at $N^{t_*}$ before the structured sieve or HB.
The endpoint $t_*$ belongs to the upper part \cite[(2.28) and the three
lines after it]{PaperII}. The coefficient is $-1$.

Stage one. For a prefix $p_1\in\mathcal P_N$, dropping the condition
$(m,Nb)=1$ gives $S(p_1,N^\alpha)\le S_{p_1}(N^\alpha)$ in the notation of
Paper~II, (5.9) \cite[the sentence after (5.9)]{PaperII}. Prefixes
$p_1\mid N$ are removed at their first failure. This is the error-table
entry ``Exceptional source factors'', of size
$O_\sigma(L^3/N+N^{-\alpha/2+o(1)})$ \cite[\S6.3]{PaperII}. A prime equal
to $N^{t_*}$ or $N^{\tau_S}$ divides $N$ \cite[after (5.20)]{PaperII}. With
the coefficient $-1$ this gives the first inequality.

Stage two. Paper~II applies the upper sieve (5.4) to each $\mathcal A_b$,
$b=p_1$. Its parameters are $X_{\mathcal A}=\operatorname{Li}(X)/\varphi(b)$, level
$Q=N^{1/2}L^{-B}$ ($\kappa=0$), $D_b=Q/b$ and $z_b=N^\alpha$
\cite[(5.12)]{PaperII}. For $t=\log_Np_1\in[t_*,\tau_S)$ one has
$a_0(t)=1/2-t\ge55/224$, $t_*-\alpha=5368241/32172000>0$, and
\[
 \frac{a_0}{\alpha}\in\Bigl[\frac{12639}{4480},\frac{7917759}{2800000}\Bigr]
 \subset(2,3)
\]
\cite[table after (5.12), row T5]{PaperII}. Hence eventually
$s_b=(a_0-\delta_N)/\alpha\in(2,3)$, where $F(s)=2e^\gamma/s$
\cite[between (5.4) and (5.5)]{PaperII}. By (5.14) the remainders are $o(U_C)$. The main
terms are (5.15), and (5.16) gives
$T_5\le4(1-\varepsilon)U_CJ_5(0)+o(U_C)$ with the vanishing logarithmic
loss (5.21) \cite[(5.14)--(5.16), (5.21) and the sentence ``S5 keeps its
vanishing logarithmic loss'']{PaperII}. Since $W=4(1-\varepsilon)U_C$, this
is the second inequality. Its error is the error-table entry ``S5 BV
remainder'' plus the prime-measure and Euler conversion of that row
\cite[\S6.3]{PaperII}.

Stage three. By Paper~II, (A.14)--(A.14a), $4J_5(0)$ lies in the ball of
centre $c_5$ and radius $8.84\cdot10^{-57}$, so $J_5(0)\le\overline J_5$.
Paper~II also gives the independent rational enclosure
$P_5<4J_5(0)<P_5+R_5$ inside the same ball \cite[after (A.14a)]{PaperII}.

The chain uses $\kappa=0$, so the increment $J_5(\kappa)-J_5(0)$ does not
occur. Paper~II's allowance $E_{5,\kappa}=4\cdot10^{-8}$ is a summand of
$g_{\mathrm{BV}}=0.00000575$ in \eqref{eq:retB-fixed-cost} and stays there
once. Paper~II, Appendix~A, records the hypothetical increment and its
once-booked decomposition. They are in the paragraph opening
\emph{For comparison, if S5-top used $\kappa=10^{-6}$} and in the
BV-allowance line right after it
\cite[Appendix~A, comparison paragraph and following decomposition]{PaperII}.
Equation (A.14a) gives only the nominal value at $\kappa=0$.
\end{proof}

The chain $\mathrm{T5}$ does not meet the changes of this section. Its
three stages use no sieve level $\ell_c$ and no saved profile. They use no
Class II, packing or Class I term. Paper~II's prepared Class II, packing and Class I
contributions of this row are zero \cite[\S2.4, paragraph ``S5-top and S6
use BV separately'']{PaperII}. The lift below changes the kernels on $I_4$,
$I_5=[1/4,t_*)$, $I_7$ and $\mathcal D_{10}$ only. The interval $I_5$ is disjoint
from $[t_*,\tau_S)$. The joint shift changes the $G_1$ cone only. Hence
Lemma~\ref{lem:retB-S5top} holds unchanged after both changes.

Keep the whole fixed deduction of Paper~II:
\begin{equation}\label{eq:retB-fixed-cost}
\begin{aligned}
 F_{\mathrm{net}}={}&0.000006484294227956901+0.00000575
                   +0.00002853328\\
                  &+4.98\cdot10^{-54}+2\cdot10^{-24}+10^{-390}\\
 &<0.000040767574227956903.
\end{aligned}
\end{equation}
These are Paper~II's six deductions \cite[(6.6)]{PaperII}.
The first remains paid although its source is excluded.
The joint $G_1$ reconciliation also remains paid.
Every adverse row difference has an absolute bound.
Keep their whole vanishing envelope, $E_{\mathrm{van}}$.
The row estimates required here are
\begin{equation}\label{eq:retB-base-errors}
 \sum_{j\in\mathcal I}\sum_{k=1}^{10}
       (Q_{j,k-1}-Q_{j,k})
 \ge-\bigl(F_{\mathrm{net}}-\delta_{G1}\bigr)
       -E_{\mathrm{van}}(N;\sigma),
 \qquad E_{\mathrm{van}}(N;\sigma)\longrightarrow0.
\end{equation}
The $\delta_{G1}$ part is already in the endpoint sum, so the remaining
chain deduction is $F_{\mathrm{net}}-\delta_{G1}$. Their total is exactly
$F_{\mathrm{net}}$.
The fixed AA37 share is paid under the admissible condition
$0<\sigma\le4g_{\mathrm{AA}}/(C_EW_*)$ of Paper~II, (1.27)
\cite{PaperII}. See Paper~II, \S6.5. This condition is part of the
fixed-parameter choice made before the limit.
The analytic inputs are Paper~II, (6.4b)--(6.5),
and its direct and $G_1$ estimates
\cite[(4.66)--(4.69), (6.4b)--(6.5)]{PaperII}.
Their logarithmic parameters are
\begin{equation}\label{eq:retB-precision}
 K_0=8042,\qquad A'=8047,\qquad A_{\mathrm{work}}=8051,
 \qquad K_0+2-A'=-3.
\end{equation}
These are Paper~II, (4.66), for the original coefficient families.
This section leaves them unchanged.

\begin{lemma}[Inherited chain errors]\label{lem:retB-base-errors}
The bound \eqref{eq:retB-base-errors} holds, and $E_{\mathrm{van}}(N;\sigma)\to0$
as $N\to\infty$ at fixed $\sigma$. Each row or separately listed stage
difference used below has an allocated nonnegative share in Paper~II's
error table. Denote it by $\mathcal E_j$. Paper~II, \S6, charges these shares once.
An entry may be grouped across rows, so the lemma claims nothing for
arbitrary subfamilies and does not claim pairwise disjoint entries. The only
omission bounds used are those for the whole chains or stage differences
named in the proof below.
\end{lemma}

\begin{proof}
The signed-assembly theorem of Paper~II proves
$D_{\mathrm{src}}\ge B_{\mathrm{main}}-F_{\mathrm{pers}}-E_{\mathrm{van}}$
with $E_{\mathrm{van}}\to0$ for the source (6.1)
\cite[\S6.1, (6.1)--(6.2), (6.4)--(6.4d)]{PaperII}. Its proof bounds every
adverse stage difference by an allocated share of Paper~II's error table,
\S6.3. Each entry is a nonnegative bound or a proved zero; some entries are
grouped across rows. Paper~II assigns each tuple to its first failed
condition and charges each difference or part once below
\cite[\S6.3]{PaperII}. The fixed entries are the six deductions of
\eqref{eq:retB-fixed-cost}, that is $F_{\mathrm{pers}}=F_{\mathrm{net}}$
\cite[(6.6)]{PaperII}. The others tend to zero at fixed $\sigma$, and their
sum is $E_{\mathrm{van}}$. The rows of $\mathcal I$ are rows of (6.1).
The $S_6$ level increment is still deducted when that chain is omitted.
The $\delta_{G1}$ part is already
in \eqref{eq:retB-base-endpoints}. This gives \eqref{eq:retB-base-errors}.
For $\mathrm{T5}$ the differences and entries are those of
Lemma~\ref{lem:retB-S5top}.
Separately listed entries cover the omissions used in this chapter.
The whole $S_6$ and $S_{\mathrm{sw}}$ chains use the S6 BV remainder,
S6 level increment and $E_{S_{\mathrm{sw}}}$ entries. The whole $G_1$
chain uses the \S4.6 table and (4.85) of Paper~II
\cite[\S4.6, (4.85)]{PaperII}. The stage-nine strip for the retained
part of $S_{10}$ uses the S10 selector entry. These bounds apply
to the named whole chains or stage difference after allocating their shares.
No arbitrary omitted subfamily is asserted. For the direct rows, the stage-nine entry
``Profile comparison'' is zero adverse. The stage-ten entry ``Directed
enclosure'' is zero additional \cite[\S6.3]{PaperII}. For the retained part
of $S_{10}$, the strip between the retained cells and $\mathcal D_{10}$ is the
entry ``S10 selector'', $7.160\tau$. The term $g_\tau=2\cdot10^{-24}$ replaces it
\cite[\S6.3, rows ``S10 selector'' and $g_\tau$]{PaperII}.

The entries rest on the following inputs, all used by Paper~II on these
same sources.
The linear-sieve main terms and remainders are Paper~I, Theorem~2.6,
(2.30)--(2.33b). This is the structured form of Iwaniec's Theorem~4
\cite[Theorem~4, p.~311]{Iwaniec1980Remainder}, applied with $D=L_c$,
$\epsilon_s=\sigma$ and $\delta_s=\tau$ \cite[\S6.2]{PaperII}.
Before the exterior sign, applied once, each block remainder splits as
\[
 \frac{C_c}{M_c}\langle\gamma_c,\Delta_P\rangle+NC_{\mathrm{sf},c}
 -\langle\Xi_{U,c},\Delta_P\rangle+\langle\Xi_{A,c},\Delta_0-\Delta_P\rangle,
\]
with $\Xi_{A,c}=O_c\mathbf 1_{q\ \mathrm{not\ squarefree}}$
\cite[(6.4a)--(6.4b), (4.23)]{PaperII}.
The term $NC_{\mathrm{sf}}$ vanishes on the direct rows, where
$\Delta_0=\Delta_P$. On $G_1$ it is $O_\sigma(N^{-\alpha}L^5)$ by Paper~II,
(4.64)--(4.65). The terms $\Xi_U$ and $\Xi_A$ are bounded jointly, once, by
$232\Lambda_\sigma C_{\mathrm{ret},\sigma} N^{-\ell_0\sigma^2}(1+L)^4$, with
$\ell_0=264897/2060000$ \cite[(B.1)--(B.7)]{PaperII}.
Here $C_{\mathrm{ret},\sigma}$ is the retained-side combinatorial constant;
it depends on the fixed $\sigma$ and not on $N$.
The persistent part $C_E\sigma W_{\mathrm{AA}}/4$ of Paper~II, (B.8), lies in
$g_{\mathrm{AA}}$; only its secondary part belongs to $E_{\mathrm{van}}$.
The Class II estimate (6.5) of Paper~II applies Paper~I, Theorem~3.2, (3.13),
with the precisions \eqref{eq:retB-precision}. The spectral input of that
theorem is Paper~I, Theorem~2.8; see Fact~\ref{fact:spectral-implication}.
In the large-norm case of Paper~I, (3.48), its mean-square input is
Lemma~\ref{hyp:ra-ms-transfer}(ii), in place of \cite[Theorem~0(a)]{BFI86}.
Lemma~\ref{hyp:ra-ms-transfer}(ii) applies under the Paper~I
hypotheses and selected parameter box stated there.
It uses neither the optional six-endpoint lemma nor any broader BFI86 scope.
Its premises are checked for these sources as follows. In Paper~I, conditions (3.1), (3.2), (3.6), (3.8) with
the length condition, and (3.11) display its hypotheses (H1)--(H5). Paper~II verifies them for these same sources, coefficients,
boxes and scales \cite[\S6.2, hypothesis table, row ``C1--C5, LEN'']{PaperII}.
Paper~I proves condition (3.25) for every selected factor of this form.
The proof uses the prime Siegel--Walfisz theorem (2.21) \cite[(3.25) and its
proof]{PaperI}. Conditions (3.26)--(3.27) are the definition of $\lambda_r$
and the Cauchy--Schwarz step of that proof. They hold under (H1)--(H5), with
$|\lambda_r|\le\tau(r)$ from the roughness in (H2) \cite[(3.26)--(3.27)]{PaperI}.
Stages one to eight are unchanged, so these are Paper~II's own checks. The prime-measure and Euler
conversions are the prime number theorem, Mertens' theorem and Paper~II,
(4.76)--(4.78).
The profile lift and the shifted joint profile below act at stages nine
and ten only. They change no source, cell, level, mask, box, norm or
precision at stages one to eight. Hence $E_{\mathrm{van}}$ is Paper~II's own
envelope for these rows.
\end{proof}

Paper~II's rational floor satisfies
\begin{equation}\label{eq:retB-paperII-floor}
 B_{\mathrm{print}}-F_{\mathrm{net}}
 >D_{\mathrm{II}}
 =\frac{45849916338009170593}{25000000000000000000};
\end{equation}
see \cite[(6.3a), (6.6)]{PaperII}.
The selected S2 endpoint used in this floor is $d_{T8}$.
Paper~II's tail reserve is
\begin{equation}\label{eq:retB-T8-reserve}
 \delta_{T8}=\frac1{20000},\qquad
 d_{T8}=\frac{P_--(2/e^\gamma)V_{\mathrm{full}}
                 -(2/e^\gamma)J_++C_-}{4}-\delta_{T8}.
\end{equation}
Thus $\delta_{T8}$ is deducted once inside $d_{T8}$, as in
Paper~II, (A.7)--(A.8). It is not a second deduction from
$D_{\mathrm{II}}$. The separate fixed term $275009680\tau$ remains
outside $d_{T8}$ in the S2 estimate \cite[(3.22)--(3.23), (A.7)]{PaperII}.
Combining the preceding estimates gives
\begin{equation}\label{eq:retB-base}
 \frac{R_{\mathrm{fix}}+S_2}{W}
 \ge D_{\mathrm{II}}+\overline J_6+2\overline H
                         -E_{\mathrm{van}}(N;\sigma).
\end{equation}
This deduction uses the individual signed chains.
Its analytic input is \eqref{eq:retB-base-errors}.

\subsection{Six improved endpoints}

Let
\begin{equation}\label{eq:retB-six-index}
 \mathcal J=\{1,2,4,5,7,10\}.
\end{equation}
For $j\in\mathcal J$, let $P_j$ be its Paper~II profile integral.
The indices $5$ and $10$ keep the preceding conventions.
Let $O_j$ and $N_j$ be directed lower endpoints.
The required numerical comparisons are
\begin{equation}\label{eq:retB-six-comparison}
 P_j\ge N_j\ge O_j,\qquad
 P_j-O_j=(N_j-O_j)+(P_j-N_j).
\end{equation}
The old and improved directed endpoints are the following exact terminating
decimals. They are given in the normalized D
contributions used by the signed sum. For S4, S5-low and S7 the sign $-1/4$ is included; for S10
the sign $+1/4$ is included.
\begin{center}
\small
\renewcommand{\arraystretch}{1.08}
\begin{tabular}{@{}lp{0.68\linewidth}@{}}
\hline
Endpoint source & Directed contribution in D units\\
\hline
S1, \cite[(A.4), (A.5a)]{PaperII} & $O_1=9.67653874894124412908238749318694090949698021569432709634$\\
 & $N_1=9.676557035420731097796376157212899872020518126053458552042$\\
S2, \cite[(A.8), (A.8a)]{PaperII} & $O_2=1.6195813569890835255913340433210572922007671408165$\\
 & $N_2=1.62063469556829$\\
S4, \cite[Appendix~A.2--A.6]{PaperII} & $O_4=-0.8342438370724665$\\
 & $N_4=-0.8341817006889064605174981431060625$\\
S5-low, \cite[Appendix~A.2--A.6]{PaperII} & $O_5=-0.05972210829776175$\\
 & $N_5=-0.05971885506770766576732140865035175$\\
S7, \cite[Appendix~A.2--A.6]{PaperII} & $O_7=-0.574850630119163$\\
 & $N_7=-0.574815106890481420929612425345365$\\
S10-Ret, \cite[Appendix~A.2--A.6]{PaperII} & $O_{10}=0.214711863944485$\\
 & $N_{10}=0.21559427153309934685978873940373255$\\
\hline
\end{tabular}
\end{center}
The rational endpoint differences sum to
\begin{equation}\label{eq:retB-six-gain}
\begin{aligned}
 g_6&=\sum_{j\in\mathcal J}(N_j-O_j)\\
 &=0.002054945489603492768011383006854970322770769542631455702.
\end{aligned}
\end{equation}
Exact rational addition of the six entries above gives the displayed
sum, and each entry satisfies $N_j\ge O_j$.
The values $O_j$ are Paper~II's directed endpoints, and $P_j\ge O_j$ is
proved there \cite[(A.4), (A.5a), (A.8), (A.8a), Appendix~A.2--A.6]{PaperII}.
The values $N_j$ are lower bounds read from Paper~II's printed independent
enclosures. The following optional lemma states them.
\begin{lemma}[Six directed endpoints from Paper II; optional]
\label{lem:retB-six-endpoints}
For each $j\in\mathcal J=\{1,2,4,5,7,10\}$, the Paper~II profile integral
$P_j$, with its signed normalized coefficient, satisfies $P_j\ge N_j$,
with $N_j$ as tabulated above.
\end{lemma}
\begin{proof}
For $j=1$, Paper~II, (A.5a), gives the printed $3/4$ ball; its lower endpoint
is $N_1$. For $j=2$, Paper~II, (A.8a), gives the printed interval whose
lower endpoint is $N_2$. Paper~II states that this interval covers the continuous functional and the infinite
tail. The $d_{T8}$ deduction in (6.6a) is released once in $N_2-O_2$.
The term $F_{\mathrm{net}}$ in (6.6b) contains no $\delta_{T8}$
\cite[(6.6a)--(6.6b)]{PaperII}. For $j=4,5,7$, the Appendix~A.3 balls
include the adverse $\tau=10^{-90}$ shifts. Since the profile function is
nonincreasing, the shifted raw integrals are at least the $\tau$-free ones.
Their upper endpoints with the negative coefficient therefore give
$P_j\ge N_j$. For $j=10$, the Appendix~A.3 $S_{10}$-Ret ball encloses the
$\tau$-free actual-coordinate profile on $\mathcal D_{10}$. Its lower
endpoint gives $P_{10}\ge N_{10}$. The separate $D_*\tau$ strip is paid in
the $S_{10}$ selector estimate. These are the printed enclosures in Paper~II,
(A.5a), (A.8a), and Appendix~A.3. Its numerical-certificate section records
the adverse-shift and actual-coordinate conventions
\cite[(A.5a), (A.8a), Appendix~A.3]{PaperII}. The remaining directed
endpoints $O_j$ are unchanged.
\end{proof}
With $O_j$ in place of $N_j$, \eqref{eq:retB-six-comparison} is
Paper~II's own endpoint comparison and $g_6$ is replaced by $0$.
Only the endpoint changes in \eqref{eq:retB-telescope}.
All earlier signed differences remain unchanged.
Thus \eqref{eq:retB-six-comparison} adds $g_6$ once.
The $S_5$ upper part remains separate.
The joint $G_1$ endpoint is absent from $\mathcal J$.

\subsection{Four stronger profile comparisons}

Use the sieve functions of Paper~I
\cite[(2.23), immediately before Theorem~2.6]{PaperI}. Put
\begin{equation}\label{eq:retB-hat-functions}
 \widehat F=\frac{F}{2e^\gamma},\qquad
 \widehat f=\frac{f}{2e^\gamma},\qquad
 K=\frac{57}{64}-u.
\end{equation}
Define
\begin{equation}\label{eq:retB-H6-domains}
\begin{aligned}
 s(x)&=\frac{K-5x/2}{\alpha},\qquad
 q(x,y)=\frac{K-5x/2-y}{\alpha},\\
 B_0(x)&=\min\{\beta,\max\{\alpha,K-11\alpha/5-5x/2\}\}.
\end{aligned}
\end{equation}
The domains are
\begin{equation}\label{eq:retB-H6-domain-sets}
\begin{gathered}
 I_4=[\mu_*,1/4),\quad I_5=[1/4,7109/28000),\quad
 I_7=[\mu_*,\lambda),\\
 \mathcal D_{10}=\{(x,y):\mu_*\le x<\lambda,\ \alpha\le y<B_0(x)\}.
\end{gathered}
\end{equation}
These are Paper~II's domains
\cite[(2.27), (2.29), (2.31), (2.38)--(2.40)]{PaperII}.
Set
\begin{equation}\label{eq:retB-H6-gains}
\begin{aligned}
 \Delta_j&=\int_{I_j}
  \frac{\widehat F(s(x))-\widehat F(s(x)+h/\alpha)}{\alpha x}\,dx
                      &&(j=4,5,7),\\
 \Delta_{10}&=\iint_{\mathcal D_{10}}
  \frac{\widehat f(q(x,y)+h/\alpha)-\widehat f(q(x,y))}
       {\alpha xy}\,dx\,dy.
\end{aligned}
\end{equation}
The three upper rows carry a negative sign.
The lower $S_{10}$ row carries a positive sign.
Their original coefficients are included in these gains,
so there is no further factor $1/4$.

Paper~II's actual levels give
\begin{equation}\label{eq:retB-H6-margins}
\begin{aligned}
 \ell_c&\ge\alpha s(x)+h+\frac1{10000}-3\tau
                   &&(j=4,5,7),\\
 \ell_c&\ge\alpha q(x,y)+h+\frac1{10000}
                         -\left(\frac{5\lambda}{2}+\beta\right)\tau
                   &&(j=10).
\end{aligned}
\end{equation}
See \cite[\S2.4, proof of (2.26)--(2.31);
\S2.5, proof of (2.40)]{PaperII}.
The exact split behind the displayed reserve is
\begin{equation}\label{eq:retB-level-buffer}
 u-\eta_L=\frac{19}{10000}
 =\frac9{5000}+\frac1{10000}=h_G+\frac1{10000}.
\end{equation}
After the $h_G$ lift, Paper~II's S4, S5-low (including its B branch)
and S7 comparisons leave $1/10000-3\tau$. Its S10-Ret comparison leaves
$1/10000-\tau(5\lambda/2+\beta)$.
Both are positive for $\tau\le10^{-90}$.
The G1-3 shift leaves
$g_s=1/10000-(225/256)\tau>2\mu_*\tau$ and the $\mu_*$ crossing
leaves $1/10000-(75/256)\tau>0$, by \cite[(4.72)]{PaperII}.
These are pointwise profile-numerator inequalities. They do not change
the actual level $\ell_c$, the level $D_c$, or the masks, boxes,
Class-II and AA37 inputs. The other level conditions of Paper~II
remain pointwise and unmodified. They include the G1 guard, S1--S3,
S2/S3/S8/S9, and the whole-K margin. See Paper~II, (2.21), (2.40),
(4.72), and (4.19).
Decreasing $F$ gives the first three comparisons.
Increasing $f$ gives the fourth comparison.
These operations leave the actual levels unchanged.
The H6 lift and the M5 shift act only at stages 9--10, after Class II
at stage 7. They change no Class-II, AA37, box, mask or precision input:
$\ell_c=\Theta_c-\eta_L-\sum_i t_i$ contains no saved profile, and
\eqref{eq:retB-precision} remains unchanged.

Proposition~\ref{prop:app-h6-four-profiles} gives finite integral
enclosures. With the directed calculation of the computational appendix,
they prove the following. Its exact rational interval is
\par\noindent
$[0.005720047094327988714222991186379229149381887590042773120871489672965022499
\mathbin{\pm}7.09\cdot10^{-76}]$.
This is a ball for the sum of the four rowwise lower bounds, not for
$\sum_j\Delta_j$.
We use the terminating credit $572004/10^8$. It lies below
that exact lower endpoint and leaves the existing $10^{-20}$ allowance intact.
This allowance is conservative; it is not an additional gain in the final constant.
\begin{equation}\label{eq:retB-H6-finite}
 \frac{572004}{10^8}
 <\Delta_4+\Delta_5+\Delta_7+\Delta_{10}
 <\frac{606486}{10^8},\qquad
 \frac{572004}{10^8}-10^{-20}>\frac{5719}{10^6}.
\end{equation}
The unchanged $S_{10}$ strip also satisfies
\begin{equation}\label{eq:retB-S10-strip}
 \frac{10987/10000}{4\alpha^2}
 \left(\frac{5\lambda}{2}+\beta\right)
 \frac{\lambda-\mu_*}{\mu_*}
 =\frac{21137254516197687}{4120000000000000}
 <5.137<7.160.
\end{equation}
The lifted $S_{10}$ argument satisfies
$q+h/\alpha\le11519839/3200000+h/\alpha<4$.
Thus $f\le f(4)=2e^\gamma\log3/4$ throughout the strip,
as used in this density bound.
These are finite bounds, including their stated domains.
The retained cells also leave a reserve. Put
\begin{equation}\label{eq:retB-S10-surplus}
\begin{gathered}
 R_{10}=\Bigl[\frac15,\frac{21}{100}\Bigr]\times\Bigl[\frac9{100},\frac1{10}\Bigr],
 \qquad \delta_{10}=\frac1{20000\alpha},\\
 c_{10}=\frac1\alpha\cdot\frac{7\delta_{10}}{480}\cdot\frac1{210}
 =\frac{146689}{320000000000},\qquad
 c_{10}-\frac{21137254516197687}{4120000000000000}\,\tau>10^{-7}.
\end{gathered}
\end{equation}
Lemma~\ref{lem:retB-lifted-direct} uses \eqref{eq:retB-S10-strip} and
\eqref{eq:retB-S10-surplus}, so its comparison over the whole of
$\mathcal D_{10}$ needs no deduction. Paper~II's $S_{10}$ selector entry
$7.160\tau$ is replaced in Paper~II by $g_\tau=2\cdot10^{-24}$, a summand of
\eqref{eq:retB-fixed-cost}. The fixed $g_\tau$ amount is deducted once.
The retained reserve above eliminates the $S_{10}$ selector deduction in
this comparison. The inherited $G_1$ peel and its label replacement are
jointly covered by that same summand.

We next define the finite profile arrays.
Let $X_0=[u,1]$. On this compact interval, use the zero extension of the
profile outside its displayed domain, and put
\begin{equation}\label{eq:retB-measure}
 m_N=\sum_{\substack{p\in\mathbb P,\ p>2,\ p\nmid N\\N^u\le p\le N}}
       \frac{\delta_{\log_Np}}{p-1},\qquad m(dt)=\frac{dt}{t}.
\end{equation}
The symbol $\delta_x$ denotes unit mass at $x$.
Define
\begin{equation}\label{eq:retB-Euler}
 V_N(z)=\prod_{\substack{p<z,\ p\in\mathbb P\\p\nmid N}}
           \left(1-\frac1{p-1}\right),
\end{equation}
and
\begin{equation}\label{eq:retB-direct-normalization}
\begin{aligned}
 X_D&=\sum_{v\in\mathbb P}
        \mathbf1_{N/L^{30}<v<(1-\varepsilon)N},\\
 r_{D,N}&=\frac{2e^\gamma\alpha X_DV_N(N^\alpha)}{W}.
\end{aligned}
\end{equation}
For $j=4,5,7$, set
\begin{equation}\label{eq:retB-direct-base-kernels}
\begin{aligned}
 k_j^0(x)&=-\frac{\mathbf1_{I_j}(x)}{\alpha}\widehat F(s(x)),\\
 k_j^h(x)&=-\frac{\mathbf1_{I_j}(x)}{\alpha}
                       \widehat F(s(x)+h/\alpha).
\end{aligned}
\end{equation}
For $j=10$, set
\begin{equation}\label{eq:retB-pair-base-kernels}
\begin{aligned}
 k_{10}^0(x,y)&=\frac{\mathbf1_{\mathcal D_{10}}}{\alpha}\widehat f(q(x,y)),\\
 k_{10}^h(x,y)&=\frac{\mathbf1_{\mathcal D_{10}}}{\alpha}
                            \widehat f(q(x,y)+h/\alpha).
\end{aligned}
\end{equation}
Put $e_j=1$ for $j=4,5,7$, and $e_{10}=2$.
The original and lifted arrays are
\begin{equation}\label{eq:retB-direct-arrays}
\begin{aligned}
 P_j&=\int k_j^0\,d(m^{e_j}),&
 S_{j,N}^0&=r_{D,N}\int k_j^0\,d(m_N^{e_j}),\\
 S_{j,N}^h&=r_{D,N}\int k_j^h\,d(m_N^{e_j}),&
 B_{j,N}&=S_{j,N}^h-S_{j,N}^0.
\end{aligned}
\end{equation}
The powers denote product measures.
These formulas use $P_j$ from \eqref{eq:retB-six-comparison}.
Write $d_j=k_j^h-k_j^0$. For $j=4,5,7$, and then $j=10$,
\begin{equation}\label{eq:retB-direct-kernels}
\begin{aligned}
 d_j(x)&=\frac{\mathbf1_{I_j}(x)}{\alpha}
  [\widehat F(s(x))-\widehat F(s(x)+h/\alpha)],\\
 d_{10}(x,y)&=\frac{\mathbf1_{\mathcal D_{10}}}{\alpha}
  [\widehat f(q(x,y)+h/\alpha)-\widehat f(q(x,y))].
\end{aligned}
\end{equation}
Thus $\Delta_j=\int d_j\,d(m^{e_j})$.
Define the two conversion errors:
\begin{equation}\label{eq:retB-direct-Lambda}
 \Lambda^0_{j,N}=S^0_{j,N}-P_j,\qquad
 \Lambda_{j,N}=B_{j,N}-\Delta_j\qquad(j\in\mathcal H).
\end{equation}

\begin{lemma}[Finite profile shift]\label{lem:retB-finite-profile}
For the arrays above,
\begin{equation}\label{eq:retB-H6-use}
 S_{j,N}^h=P_j+\Lambda^0_{j,N}+\Delta_j+\Lambda_{j,N}
 \ge P_j+\Delta_j-|\Lambda^0_{j,N}|-|\Lambda_{j,N}|.
\end{equation}
\end{lemma}
\begin{proof}
Substitute the two definitions in \eqref{eq:retB-direct-Lambda}.
\end{proof}
Both absolute conversion terms appear as separate summands of $E'$.
We do not assume that Paper~II's original envelope contains the
saved-array error. Its stage 8 converts the actual-level array before
stage 9 domination \cite[\S6.1, stage table]{PaperII}.
The arithmetic application uses the preceding level domination.

\paragraph{Restoration.}
The four rows of $\mathcal H$ are Paper~II's rows $S_4$, the lower part
of $S_5$, $S_7$ and the retained part of $S_{10}$. Their stage-two
source-restoration shares are assigned once to the source parents $S_4$,
$S_5$, $S_7$ and $S_{10}$. The literal $S_5$ split at $7109/28000$
is made before the structured sieve or Heath--Brown expansion, as above.
Both parts retain the parent $S_5$ for this restoration accounting
\cite[(2.28), \S6.1, stage table and the list of S10, S5 and S4/S7 endpoints]{PaperII}.
There is no separate restoration for the lower part of $S_5$. The upper part
of $S_5$ is the outer chain $\mathrm{T5}$ of Lemma~\ref{lem:retB-S5top},
which the lift does not change. Each tuple of an actual source is assigned
to one entry by Paper~II's first-failure rule \cite[(6.4d) and the two
sentences after it]{PaperII}. This section changes no source and no stage
from one to eight, so the assignment is Paper~II's. Appendix~\ref{subsec:app-retB-rows} lists the finite
correspondence between these rows and Paper~II's error entries. It is an index of that
assignment, and each entry occurs once.

\begin{lemma}[Lifted direct comparison]\label{lem:retB-lifted-direct}
Let $j\in\mathcal H$, and let $Q_{j,8}$ be Paper~II's stage-eight quantity
for this row. It is the continuum mass of the actual-level main term after the
prime-measure and Euler conversion. There each actual cell carries its
actual relative level $\ell_c$. Put $Q^h_{j,9}=P_j+\Delta_j=\int
k_j^h\,d(m^{e_j})$. Then $Q_{j,8}\ge Q^h_{j,9}$.
\end{lemma}

\begin{proof}
Stages one to eight are Paper~II's. Paper~II passes from $Q_{j,8}$ to
$Q_{j,9}=P_j$ by pointwise domination of the actual-level integrand by the
specified profile. The cells and the continuum measure are the same
\cite[\S6.1, stage table; \S2.4, after (2.31); \S2.5, after (2.40)]{PaperII}.
We repeat that domination with the lifted profile.
For each B-labelled parent contributing to $S_4$, $S_7$ or $S_5$-low,
we first rederive the buffer used below. Write
\[
 m(t)=\min\left\{\frac{1+t}{2},\frac{57}{64}-\frac{3\rho t}{2}\right\},
 \qquad t_{\rm dd}=\frac{25}{32(1+3\rho)}.
\]
Paper~II gives $\Theta_c\ge m(t_1)$. The branch comparison is the exact
identity
\[
 \frac{1+t_1}{2}
 -\left(\frac{57}{64}-\frac{3\rho t_1}{2}\right)
 =\frac{1+3\rho}{2}(t_1-t_{\rm dd}).
\]
Thus the decreasing branch is active exactly when $t_1>t_{\rm dd}$.
Equality is assigned to A, where both bounds agree. This active branch test
is applied to every B-labelled parent. The branch is never inferred from the
largest completed coordinate alone. Paper~II keeps the inherited label through
clipping. In S4 and S7 the B parents are separate intersections of their
respective row domains with $t_1>t_{\rm dd}$. No common parent interval is
asserted. S5-low has $1/4\le r<t_*=7109/28000$ and
$t_1\le r<\rho t_1$, so $t_1>1/(4\rho)>t_{\rm dd}$ because $\rho<8$.
It is therefore entirely B. Let a parent in S4 or S7 be A-labelled.
Then $t_1\le t_{\rm dd}$, and Paper~II's first branch in (2.17f)
\cite[(2.17f)]{PaperII} gives
$\Theta_c\ge(1+t_1)/2$. Take an actual exponent $r\in[t_1,\rho t_1)$.
The S4 and S7 row domains are clipped by $r\ge\mu_*$
\cite[\S2.1, fixed-cutoff rows]{PaperII}, and the retained bin gives
$r\ge t_1$.
The residual is therefore
\[
 2(t_1-\mu_*)+\frac52(r-t_1)
 =2(r-\mu_*)+\frac12(r-t_1)\ge0.
\]
\[
 \ell_c-\left({57\over64}-u-{5r\over2}\right)
 \ge u-\eta_L-{25\over64}+2t_1+{5\over2}(r-t_1)
 = {19\over10000}+2(t_1-\mu_*)+{5\over2}(r-t_1)
 \ge {19\over10000}.
\]
Here $\mu_*=25/128$ is the retained cutoff parameter; the symbol $\mu(n)$
elsewhere denotes the Möbius function.
Thus the A pieces also have the required $19/10000-3\tau$ buffer. This is a row-local calculation.
It does not merge the distinct S4 and S7 parent domains. For any B parent,
Paper~II's floor gives
\begin{equation}\label{eq:retB-B-branch-buffer}
 \Theta_c\ge {57\over64}-{3\over2}\rho t_1,\qquad
 \ell_c-\left({57\over64}-u-{5r\over2}\right)
 \ge (u-\eta_L)+{5\over2}(r-t_1)-{3\over2}\tau t_1
 \ge {19\over10000}-3\tau .
\end{equation}
The first bound is Paper~II's linear branch of the full-anatomy floor
\cite[(2.17d)--(2.17f)]{PaperII}. Using
$\ell_c=\Theta_c-\eta_L-t_1$, $r\ge t_1$, $u-\eta_L=19/10000$, and
$t_1<t_*<1$ proves the second line by direct subtraction. The B pieces of
S4 and S7 use the same row-local calculation. Their parent intervals
stay distinct.

For $j=4,5,7$ the row coefficient is negative. At a point $x$ of an actual
cell, the actual sieve argument is $\ell_c/\alpha$ and, by
\eqref{eq:retB-H6-margins},
\[
 \frac{\ell_c}\alpha\ge s(x)+\frac h\alpha+\frac{10^{-4}-3\tau}\alpha
 >s(x)+\frac h\alpha .
\]
These are Paper~II's margins $19/10000-3\tau$ for $S_4$, $S_7$ and the
$B$ branch of the lower part of $S_5$. Here $h$ has moved into the profile
\cite[\S2.4, display before ``The S4/S7 A/B levels give the same buffer'']{PaperII}.
Since $F$ is decreasing, $-\widehat F(\ell_c/\alpha)\ge-\widehat F(s(x)+h/\alpha)$.
For $j=10$ the coefficient is positive, and on the retained cells
\[
 \frac{\ell_c}\alpha\ge q(x,y)+\frac h\alpha
 +\frac{10^{-4}-(5\lambda/2+\beta)\tau}\alpha>q(x,y)+\frac h\alpha,
\]
which is Paper~II's
$\alpha(s_{\mathrm{live}}-s_{\mathrm{saved}})\ge(u-\eta_L)-\tau((5/2)\lambda+\beta)$
with $h$ moved into the profile \cite[\S2.5, proof after (2.40)]{PaperII}.
Since $f$ is nondecreasing,
$\widehat f(\ell_c/\alpha)\ge\widehat f(q(x,y)+h/\alpha)$.
Paper~II uses the same monotonicity of $F$ and $f$ for these
comparisons \cite[\S2.4]{PaperII}.
For $j=4,5,7$ the cells cover $I_j$, and integrating the pointwise
inequality against $m$ gives $Q_{j,8}\ge\int k_j^h\,dm=P_j+\Delta_j$.

For $j=10$ the retained cells need not cover $\mathcal D_{10}$. Let
$R_\tau\subset\mathcal D_{10}$ be their union, put
$a_{10}=5\lambda/2+\beta$, and write $T_h(E)=\int_Ek_{10}^h\,d(m^2)$. Then
\[
 Q_{10,8}=\frac1\alpha\sum_{c\ \mathrm{retained}}\int_c
 \widehat f(\ell_c/\alpha)\,d(m^2).
\]
\emph{Strip.} Let $(x,y)\in\mathcal D_{10}\setminus R_\tau$ lie in a cell
with labels $t_1\le x$, $t_2\le y$. The cell is dropped, so
$y\ge t_2>\operatorname{sel}(t_1)\ge B_0(t_1)-a_{10}\tau\ge B_0(x)-a_{10}\tau$.
This uses the two deficit bounds of Paper~II, \S2.5, and the fact that $B_0$ is
nonincreasing \cite[\S2.5, proof of (2.40)]{PaperII}. Also $y<B_0(x)$. On
$\mathcal D_{10}$ one has $11/5\le q$ and $q+h/\alpha<4$, so
$\widehat f(q+h/\alpha)\le\widehat f(4)=\log3/4$. With $y\ge\alpha$ and
$\int_{\mu_*}^\lambda dx/x\le(\lambda-\mu_*)/\mu_*$, this gives
$T_h(\mathcal D_{10}\setminus R_\tau)\le D_*\tau$, where $D_*$ is the
constant in \eqref{eq:retB-S10-strip}.

\emph{Reserve.} On a retained cell the margin in
\eqref{eq:retB-H6-margins} satisfies $10^{-4}-a_{10}\tau>1/20000$, so
$\ell_c/\alpha\ge q+h/\alpha+\delta_{10}$. Here
$2<q+h/\alpha<q+h/\alpha+\delta_{10}<4$. For \(2<s<4\),
$\widehat f(s)=\log(s-1)/s$ and
\[
 \widehat f'(s)=\frac{s/(s-1)-\log(s-1)}{s^2}
 >\frac{4/3-11/10}{16}=\frac7{480},
\]
since $\log3<11/10$. As $f$ is nondecreasing,
$\widehat f(\ell_c/\alpha)-\widehat f(q+h/\alpha)\ge7\delta_{10}/480$ at
every point of $R_\tau$.

\emph{A fixed retained rectangle.} Let $(x,y)\in R_{10}$, in a cell with
labels $t_1\in(x/\rho,x]$, $t_2\in(y/\rho,y]$. Then $\mu_*<t_1<\lambda$ and
$\alpha<t_2<\beta$. On $[\mu_*,21/100]$ we have $B_0=\beta$, because
$57/64-u-11\alpha/5-(5/2)(21/100)>\beta$. Hence
$\operatorname{sel}(t_1)\ge\beta-a_{10}\tau>1/10\ge t_2$, and the cell is
retained: $R_{10}\subset R_\tau$. Also
$m^2(R_{10})=\log\frac{21}{20}\log\frac{10}9\ge\frac1{21}\cdot\frac1{10}
=\frac1{210}$.

Combining the three steps with the pointwise inequality above,
\[
 Q_{10,8}\ge T_h(R_\tau)+c_{10}
 \ge T_h(\mathcal D_{10})-D_*\tau+c_{10}>T_h(\mathcal D_{10})
 =P_{10}+\Delta_{10},
\]
by \eqref{eq:retB-S10-surplus}. No deduction is made. The reserve is the
part of Paper~II's level margin left after the lift $h$. It exists on cells
of $S_{10}$ only and is used only here.
\end{proof}

The finite arrays $S^0_{j,N},S^h_{j,N}$ are not used in this inequality.
The errors $|\Lambda^0_{j,N}|$ and $|\Lambda_{j,N}|$ are nonnegative and are
kept in $E'$ so that $E'$ agrees with \eqref{eq:intro-retained-error}.
They tend to zero by Lemma~\ref{lem:retB-conversion}.

\subsection{The internal joint shift}

The joint occurrence is $2S_{11}+S_{12}+S_{13}$.
Its six-branch function is Paper~II, (4.70)
\cite[(4.70)--(4.72)]{PaperII}:
\begin{equation}\label{eq:retB-theta}
\begin{aligned}
 v_1(t)&=\min\left\{\frac{1+t}{2},\frac{32-34t}{39},
                              \frac{57}{64}-\frac{3t}{2}\right\},\\
 v(t)&=v_1(t)-u,\qquad \psi(x,y)=x\mathbf1_{x\ge y},\\
 \Theta_*(a,b,c)&=\max\{v(a),v(b),v(a+b),v(a+b+c),\\
 &\hspace{34mm}\psi(v(a+c),a+2b+c),
                  \psi(v(b+c),2a+b+c)\}.
\end{aligned}
\end{equation}
Subtract $u$ before testing each guard.
Set
\begin{equation}\label{eq:retB-G-cone}
 \mathcal C=\{(t_1,t_2,t_3):u\le t_3<t_2<t_1<\mu_*\},
 \qquad
 a_3(\mathbf t)=\frac{\Theta_*(\mathbf t)-t_1-t_2-t_3}{t_3}.
\end{equation}
Shift the output of $\Theta_*$ by $h$.
Keep its guards and actual source unchanged:
\begin{equation}\label{eq:retB-G-shift}
\begin{aligned}
 a_3^s(\mathbf t)&=a_3(\mathbf t)+\frac h{t_3},\\
 J&=\iiint_{\mathcal C}\frac{1-f(a_3(\mathbf t))}
                      {t_1t_2t_3^2}\,d\mathbf t,\\
 J_s&=\iiint_{\mathcal C}\frac{1-f(a_3^s(\mathbf t))}
                      {t_1t_2t_3^2}\,d\mathbf t.
\end{aligned}
\end{equation}
Use $f(s)=0$ for $s<2$.
Keep its true values at every larger argument.
The new margins are
\begin{equation}\label{eq:retB-G-margins}
 g_s=\frac1{10000}-\frac{225}{256}\tau>2\mu_*\tau,
 \qquad
 g_{s,\mu_*}=\frac1{10000}-\frac{75}{256}\tau>0.
\end{equation}
The second margin concerns a parent interval crossing $\mu_*$.
Both follow by subtracting $h$ in Paper~II, (4.72); see
Lemma~\ref{lem:retB-shifted-retention}.
They hold throughout $0<\tau\le10^{-90}$.
If $f(a_3^s)>0$, the retained-cell condition follows:
\begin{equation}\label{eq:retB-G-retention}
 \ell_c\ge\Theta_*+h-t_1-t_2-t_3+g_s
       >2t_3+2\mu_*\tau\ge2\rho t_3.
\end{equation}
Thus the canonical restriction still loses zero.
At fixed numerator, use the increase of $sf(s)$.
This is the same comparison as Paper~II, (4.72).

\begin{lemma}[Shifted retention]\label{lem:retB-shifted-retention}
The margins \eqref{eq:retB-G-margins} hold for $0<\tau\le10^{-90}$.
On the original joint source, every positive shifted lower contribution
lies in an original retained whole cell of Paper~II, (4.19). On such a
cell the actual lower main term dominates the shifted canonical one.
\end{lemma}

\begin{proof}
Use an original G1-3 cell with inherited labels
$t_1\ge t_2\ge t_3>0$, and order its distinct actual prefix primes as
$a_1>a_2>a_3$. Put $x_i=\log_Na_i$. Unequal parent bins enforce the
actual order, while within an equal-label group we order the primes.
Hence $t_i\le x_i<\rho t_i$, including clipped bins.
No strict inequality between the labels is required.

Let $\Theta_c$ be the original full-anatomy infimum and
$\ell_c=\Theta_c-\eta_L-\sum_it_i$. Put
$B_s(x)=\Theta_*(x)+h-\sum_ix_i$.
The function $v_1$ is $3/2$-Lipschitz. On the closed ordered parent
a coordinate changes by at most $\mu_*\tau$. A sum of at most
three coordinates changes by at most $3\mu_*\tau$.
If a saved guarded branch is active at $x$, its buffered guard has
nonnegative gap, including equality. Removing its buffer adds $u$.
Throughout the parent its unbuffered gap is therefore at least
$u-(9/2)\mu_*\tau-4\mu_*\tau>0$.
Thus it stays available in the full anatomy. Take an unguarded or
active guarded saved branch. The corresponding full branch is at least
its value at $x$ plus $u-(9/2)\mu_*\tau$. Taking their maximum and then
the parent infimum preserves this lower bound. Ties between maximizing
branches cause no loss. Therefore
\[
 \ell_c\ge B_s(x)+g_s+\sum_i(x_i-t_i),\qquad
 g_s=u-\eta_L-h-(9/2)\mu_*\tau
     =10^{-4}-(225/256)\tau.
\]
For a parent crossing $\mu_*$, the largest exponent is within
$\mu_*\tau$ of $\mu_*$, and $v_1(\mu_*)=153/256$.
Every saved branch is at most $153/256-u$, so the crossing comparison
gives instead $10^{-4}-(75/256)\tau\ge g_s$.
Clipping only shrinks the infimum domain. These arguments also apply
when two or all three labels coincide.
For $0<\tau\le10^{-90}$ we have $g_s>2\mu_*\tau>0$.

If $f(B_s(x)/x_3)>0$, then $B_s(x)>2x_3$ and hence
\[
 \ell_c>2x_3+2\mu_*\tau
       \ge2t_3+2\tau t_3=2\rho t_3.
\]
Thus the whole original cell is retained by Paper~II, (4.19).
On positive support $\ell_c>B_s(x)>0$ and $0<t_3\le x_3$ imply
\[
 \frac{\ell_c}{t_3}\ge\frac{B_s(x)}{x_3},\qquad
 f(\ell_c/t_3)\ge f(B_s(x)/x_3),\qquad
 \frac{f(\ell_c/t_3)}{t_3}
      \ge\frac{f(B_s(x)/x_3)}{x_3}.
\]
These use only nonnegativity and monotonicity of $f$. Equal labels
invalidate none of the weak inequalities.
Also $E=N^{t_3}\le a_3$, so $V_N(E)\ge V_N(a_3)$.
At zero shifted contribution the required comparison follows from
nonnegativity without any assertion that $B_s(x)>0$.

Keep the original orbit normalization. If the equal-label groups are
$G_h$, put $M_c=\prod_h|G_h|!\in\{1,2,6\}$.
For a full labelled realization $\mathbf a$, write
$\mathbf a^\downarrow$ for its decreasing rearrangement within the
equal-label groups. The shifted comparison above is evaluated at
$\mathbf x^\downarrow$ and is invariant under these group permutations.
For every nonnegative, invariant main-mass weight $w_c$, therefore,
\[
 \frac1{M_c}\sum_{\rm full}K_P(\mathbf a)w_c(\mathbf a)
       V_N(E)f(\ell_c/t_3)
 \ge
 \frac1{M_c}\sum_{\rm full}K_P(\mathbf a)w_c(\mathbf a)
       V_N(a_3^\downarrow)
       f(B_s(\mathbf x^\downarrow)/x_3^\downarrow).
\]
Each distinct ordered triple has exactly $M_c$ labelled realizations.
The prefix mask $K_P$ and its correction are kept exactly as in
Paper~II, (4.22)--(4.23). No repeated-prefix realization is reclassified.
The tail mask and all remainder coefficients are likewise unchanged.
This proves the finite lower-main-term comparison with the cell mass
divided by $M_c$. The corresponding continuum comparison uses the
preceding inequality with the $1/t_3$ and $1/x_3$ factors.
Equal-label bins are kept. Their strictly ordered actual
triples lie in the original cone. The orbit identity converts
the full-block sum to that ordered sum exactly.
Only actual coordinate diagonals have zero continuum measure.
The proof imposes no $N\to\infty$ disappearance of equal-bin regions.
\end{proof}

Put $H_{\log}=\log(\mu_*/u)$.
Paper~II's remaining terms satisfy
\begin{equation}\label{eq:retB-G-residual}
\begin{aligned}
 0\le R_G\le R_{\mathrm{up}}
 ={}&500(F(30)-1)+500H_{\log}(1-f(30))\\
    &+250H_{\log}^2(F(30)-1).
\end{aligned}
\end{equation}
Its constant terms cancel exactly
\cite[(4.73)--(4.75)]{PaperII}:
\begin{equation}\label{eq:retB-G-cancellation}
 500\left(1-H_{\log}+\frac{H_{\log}^2}{2}\right)
 -\iiint_{\mathcal C}\frac{d\mathbf t}{t_1t_2t_3^2}
 =\frac1\mu_*.
\end{equation}
The shifted upper comparison is therefore
\begin{equation}\label{eq:retB-G-profile}
 G_1^s=\frac1\mu_*+R_G+J_s,
 \qquad
 \Phi_{G_1}^s=-K_GG_1^s,
 \qquad
 K_G=\frac{2I_{11}+I_{12}+I_{13}}{2e^\gamma}.
\end{equation}
The source-mass integrals are \cite[(4.77)]{PaperII}
\begin{equation}\label{eq:retB-G-masses}
\begin{aligned}
 I_{11}={}&\int_\lambda^{1/3}
          \frac{\log((1-2t)/t)}{t(1-t)}\,dt\\
 &+\int_\lambda^{1/4}\frac{dt}{t}
        \int_t^{(1-t)/3}
          \frac{\log((1-t-2v)/v)}{v(1-t-v)}\,dv,\\
 I_{12}={}&\int_\alpha^{1/3}
          \frac{\log(2-3t)}{t(1-t)}\,dt,\\
 I_{13}={}&\int_\beta^\lambda
          \frac{\log((1-t-\lambda)/\lambda)}{t(1-t)}\,dt.
\end{aligned}
\end{equation}
We have $I_{11},I_{12},I_{13}\ge0$ and $I_{12}>0$.
Thus $K_G>0$.
The finite shifted enclosure gives
\begin{equation}\label{eq:retB-G-finite}
 0\le G_1^s\le\frac{6073714261}{10^9}
 \quad\text{(shifted mesh $n_{\mathrm{mesh}}=512$)}.
\end{equation}
Paper~II's mass enclosure gives $0\le K_G\le K_{\mathrm{hi}}$
\cite[(A.10d)]{PaperII}. The exact constants are
\begin{equation}\label{eq:retB-exact-Khi}
 K_{\mathrm{hi}}=
 \frac{3989480791070578466390443009830933753748442947478110112223}{10^{58}}.
\end{equation}
Paper~II, (4.84), gives
\begin{equation}\label{eq:retB-exact-LG1}
 L_{G1}=-
 \frac{243950428265194656254206921939947762868871361874409552095097}
 {100433627766186892221372630771322662657637687111424552206336}.
\end{equation}
Define
\begin{equation}\label{eq:retB-M5}
 M_{5,\mathrm{internal}}
   =-K_{\mathrm{hi}}\frac{6073714261}{10^9}-L_{G1}.
\end{equation}
Their exact substitution gives
\begin{equation}\label{eq:retB-M5-value}
\begin{gathered}
 0.0058749499450263575220<M_{5,\mathrm{internal}},\\
 M_{5,\mathrm{internal}}<0.0058749499450263575221.
\end{gathered}
\end{equation}
The finite enclosure and identity are proved.
They concern the shifted profile integral.
The required actual-source inequality is
\begin{equation}\label{eq:retB-G-use}
 -\frac{2S_{11}+S_{12}+S_{13}}{W}
 \ge L_{G1}+M_{5,\mathrm{internal}}
       -E^D_{G1}-|\Lambda_{G,N}|.
\end{equation}

\begin{proof}[Proof of \eqref{eq:retB-G-use}]
Paper~II proves
\[
 \frac{2S_{11}+S_{12}+S_{13}}W\le K_GG_1+E^D_{G1},\qquad
 G_1=\frac1{\mu_*}+R_G+J,
\]
from (4.73)--(4.77) \cite[(4.75), (4.78)]{PaperII}. Run the same argument
with $J$ replaced by $J_s$. The first three residual rows and the bound
$R_G\le R_{\mathrm{up}}$ do not change, and neither does the cancellation
\eqref{eq:retB-G-cancellation}. In the fourth row,
Lemma~\ref{lem:retB-shifted-retention} replaces Paper~II's comparison after
(4.72). The canonical restriction loses zero, and the actual lower main term
dominates the shifted one. The source and its weights $(2,1,1)$ do not
change, and neither do the peel (4.45)--(4.57) and the mass conversion
(4.76)--(4.78). Hence
\[
 \frac{2S_{11}+S_{12}+S_{13}}W\le K_GG_1^s+E^D_{G1}
\]
with the same $E^D_{G1}$. By Proposition~\ref{prop:app-m5-cap},
$0\le G_1^s\le6073714261/10^9$, which is \eqref{eq:retB-G-finite}, with the
residual bound supplied by \eqref{eq:retB-G-residual}. By
Paper~II, (A.10d), $0\le K_G\le K_{\mathrm{hi}}$. Both factors are
nonnegative, so by \eqref{eq:retB-M5}
\[
 -K_GG_1^s\ge-K_{\mathrm{hi}}\frac{6073714261}{10^9}=L_{G1}+M_{5,\mathrm{internal}}.
\]
This gives \eqref{eq:retB-G-use} even without $|\Lambda_{G,N}|$. That term
is nonnegative and is kept so that the error agrees with
\eqref{eq:intro-retained-error}.
\end{proof}
The existing $E^D_{G1}$ is a disjoint component of the original
$F_{\mathrm{net}}+E_{\mathrm{van}}$ allowance.
Its fixed peel $145664566.4567\tau$ is covered by $g_\tau$,
and its fixed $C_E\sigma$ share by $g_{\mathrm{AA}}$.
Its vanishing part belongs to $E_{\mathrm{van}}$.
When replacing the joint chain, carry these same components once.
The reconciliation in \eqref{eq:retB-fixed-cost} remains paid.
These components are the $G_1$ comparator term $NC_{\mathrm{sf}}$ of
Paper~II, (4.64)--(4.65), and the $G_1$ share of the joint bound. That
bound treats $\Xi_U$ and $\Xi_A$ in Paper~II, (B.1)--(B.7). Both are
attached to the original source measure and not to the profile. Hence the
shift neither creates nor removes them.

\subsection{Conversion of the finite arrays}

We give the remaining definitions before stating the lemma.
All representation variables in the following table are prime.
Their product is $e$, with $\varepsilon N<e<N$.
Let $R_{11,3}(e)$ and $R_{11,4}(e)$ count the first two forms.
Let $R_{12}(e)$ and $R_{13}(e)$ count the last two forms.
Put $R_{11}=R_{11,3}+R_{11,4}$.
These are Paper~II's representation weights
\cite[(4.4), source table preceding (4.4)]{PaperII}.

\begin{table}[ht]
\caption{The internal representation counts}\label{tab:retB-representations}
\small
\renewcommand{\arraystretch}{1.15}
\begin{tabular}{@{}llp{.64\linewidth}@{}}
\hline
Count & Product & Restrictions\\
\hline
$R_{11,3}$ & $p_1p_2p_3$ &
$N^\lambda\le p_1<N^{1/3}$, $p_1<p_2$, $p_1p_2^2<N$, $p_3\ge p_2$.\\
$R_{11,4}$ & $p_1p_2p_3p_4$ &
The same conditions on $p_1,p_2$; $p_2\le p_3\le p_4$.\\
$R_{12}$ & $p_1p_2p_3$ &
$N^\alpha\le p_1<N^{1/3}\le p_2$, $p_1p_2^2<N$, $p_3\ge p_2$.\\
$R_{13}$ & $p_1p_2p_3$ &
$N^\beta\le p_1<N^\lambda\le p_2$, $p_1p_2^2<N$, $p_1p_2p_3^2\ge N$.\\
\hline
\end{tabular}
\end{table}
For the last form the endpoint orders are distinct.
Every permitted representation keeps its multiplicity.
Define the upper cut and masses by
\begin{equation}\label{eq:retB-source-normalizations}
\begin{aligned}
 N_{\mathrm{cut}}&=N(1-L^{-333}),\\
 X_j^{\mathrm{cut}}&=\sum_eR_j(e)
           \mathbf1_{\varepsilon N<e\le N_{\mathrm{cut}}}
                &&(j=11,12,13),\\
 X_G&=2X_{11}^{\mathrm{cut}}+X_{12}^{\mathrm{cut}}+X_{13}^{\mathrm{cut}},\\
 r_{G,N}(t)&=\frac{tX_GV_N(N^t)}{WK_G}.
\end{aligned}
\end{equation}
The cut precedes the sieve decomposition
\cite[(4.16)--(4.17)]{PaperII}.
The source-mass input used below is
\begin{equation}\label{eq:retB-source-mass-limits}
\begin{aligned}
 X_D&=\left(1-\varepsilon+o(1)\right)\frac NL,\\
 X_G&=\left((1-\varepsilon)(2I_{11}+I_{12}+I_{13})+o(1)\right)
                                                   \frac NL.
\end{aligned}
\end{equation}

\begin{lemma}[Source masses]\label{lem:retB-source-mass}
The formulas \eqref{eq:retB-source-mass-limits} hold.
\end{lemma}
\begin{proof}
The sum $X_D$ counts the primes in $(N/L^{30},(1-\varepsilon)N)$, so the
prime number theorem gives $X_D=(1-\varepsilon+o(1))N/L$.
For the representation weights of Table~\ref{tab:retB-representations},
Paper~II proves $X_j=((1-\varepsilon)I_j+o(1))N/L$ for the uncut masses,
$j=11,12,13$ \cite[(4.76)--(4.77)]{PaperII}. The cut strip
$N_{\mathrm{cut}}<e<N$ carries joint mass at most $100(NL^{-333}+1)$
\cite[(4.16)--(4.17)]{PaperII}, which is $o(N/L)$. Summing with the
weights $(2,1,1)$ gives the formula for $X_G$.
\end{proof}

On $X_0^3$, define the zero-extended kernel
\begin{equation}\label{eq:retB-G-gain-kernel}
 d_G(\mathbf t)=\mathbf1_{\mathcal C}(\mathbf t)
       K_G\frac{f(a_3^s(\mathbf t))-f(a_3(\mathbf t))}{t_3}.
\end{equation}
Its arrays and integral are
\begin{equation}\label{eq:retB-G-array}
\begin{aligned}
 B_{G,N}&=\iiint r_{G,N}(t_3)d_G(\mathbf t)\,d(m_N^3)(\mathbf t),\\
 \Delta_G&=\iiint d_G\,d(m^3)=K_G(J-J_s),\\
 \Lambda_{G,N}&=B_{G,N}-\Delta_G.
\end{aligned}
\end{equation}
Equivalently,
\begin{equation}\label{eq:retB-G-prime-sum}
\begin{aligned}
 B_{G,N}&=\frac{X_G}{W}
 \sum_{\substack{p_3<p_2<p_1,\ p_i\in\mathbb P,\ p_i\nmid N\\
           u\le\log_Np_3<\log_Np_2<\log_Np_1<\mu_*}}
 \frac{V_N(p_3)[f(a_3^s(\mathbf t))-f(a_3(\mathbf t))]}
      {(p_1-1)(p_2-1)(p_3-1)},\\
 &\hspace{38mm}t_i=\log_Np_i.
\end{aligned}
\end{equation}
These definitions specify all five additional conversion errors.

\begin{lemma}[Conversion of the five profile differences]
\label{lem:retB-conversion}
Fix $0<\sigma\le10^{-10}$. For $m_N,m$ in \eqref{eq:retB-measure}, the finite
arrays \eqref{eq:retB-direct-arrays} and \eqref{eq:retB-G-array} satisfy
\begin{equation}\label{eq:retB-conversion-limit}
 \Lambda_{j,N}\longrightarrow0\quad(j\in\mathcal H),\qquad
 \Lambda_{G,N}\longrightarrow0,
\end{equation}
and $\Lambda^0_{j,N}\to0$ for $j\in\mathcal H$.
Every limit ranges over even $N$ tending to infinity; the thresholds depend
on $\sigma$.
\end{lemma}

\begin{proof}
\emph{Step 1: convergence of the harmonic measures.}
Let $u\le a<b\le1$.
Mertens' theorem and partial summation give
\begin{equation}\label{eq:retB-weak}
 m_N([a,b))=\log\frac ba+o(1).
\end{equation}
Replacing $1/(p-1)$ by $1/p$ costs $O(N^{-u})$.
There are at most $L/\log2$ prime divisors of $N$.
Each omitted divisor above $N^u$ has mass $O(N^{-u})$.
Their total cost is $O(N^{-u}L)$.
The largest endpoint atom is $O(N^{-u})$.
The total masses therefore satisfy
\begin{equation}\label{eq:retB-total-mass}
 c_N:=m_N(X_0)\longrightarrow c_0:=m(X_0)=\log(1/u)>0.
\end{equation}
Let $g$ be continuous on $X_0$. Choose a finite interval partition and a
step function on it, uniformly within a prescribed error of $g$.
The integrals of the step function converge by \eqref{eq:retB-weak}. The
integral errors are at most the uniform error times $c_N+c_0$, which is
bounded by \eqref{eq:retB-total-mass}. Hence $m_N\Rightarrow m$
\cite[Theorem~2.2 and Example~2.3]{Billingsley1999}.
For $k=2,3$, rectangle masses are products of interval masses. The same
argument with rectangular step functions gives
\begin{equation}\label{eq:retB-product-weak}
 m_N^k\Rightarrow m^k\qquad(k=1,2,3).
\end{equation}
For a coordinate diagonal $D_{ij}=\{t_i=t_j\}$,
\begin{equation}\label{eq:retB-diagonal-mass}
 m_N^k(D_{ij})\le
  \left(\max_{p\ge N^u}\frac1{p-1}\right)c_N^{k-1}
 =O(N^{-u})\qquad(k\ge2).
\end{equation}
Thus strict prime order costs nothing in the limit, and no orbit factor
enters.

\emph{Step 2: the source and Euler normalizations.}
For $p>2$,
\begin{equation}\label{eq:retB-Euler-factorization}
 1-\frac1{p-1}=
 \left(1-\frac1p\right)\left(1-\frac1{(p-1)^2}\right).
\end{equation}
Mertens' product formula and the convergent second product yield
\begin{equation}\label{eq:retB-Euler-limit}
 V_N(N^t)=\frac{2e^{-\gamma}C(N)}{tL}(1+o(1)),
 \qquad u/\rho\le t\le1,
\end{equation}
uniformly in $t$.
The prime $2$ is omitted because $N$ is even.
The missing divisors $p\mid N$ above $N^t$ contribute
$\prod(1+1/(p-2))=1+O(N^{-u/\rho}L)$. The ordinary product error is
uniform for $N^t\ge N^{u/\rho}$.
By Lemma~\ref{lem:retB-source-mass} and \eqref{eq:not-normalization},
\begin{equation}\label{eq:retB-density-limits}
 r_{D,N}\longrightarrow1,\qquad
 \sup_{u\le t\le\mu_*}|r_{G,N}(t)-1|\longrightarrow0.
\end{equation}
For example, $WK_G=2e^{-\gamma}(1-\varepsilon)
(2I_{11}+I_{12}+I_{13})C(N)N/L^2$.
This is exactly the main factor in $tX_GV_N(N^t)$.
Thus both the truncation factor and the representation multiplicities
cancel once.

\emph{Step 3: restriction and bounded kernels.}
Write $\widehat m_N=m_N/c_N$ and $\widehat m=m/c_0$.
These are probability measures, and their products converge weakly.
Let $D$ have $m^k(\partial D)=0$, and let $g\mathbf1_D$ be bounded and
$m^k$-almost everywhere continuous. Apply the mapping form of the portmanteau
theorem \cite[Theorems~2.1 and~2.7]{Billingsley1999} to
$g\mathbf1_D$ and the probability measures. Then multiply by
$c_N^k\to c_0^k$. This gives
\begin{equation}\label{eq:retB-portmanteau}
 \int_Dg\,d(m_N^k)\longrightarrow\int_Dg\,d(m^k).
\end{equation}
The intervals $I_j$ have finitely many boundary points.
The boundary of $\mathcal D_{10}$ is piecewise affine, since $B_0$ is
continuous and piecewise affine. Both have zero $m^{e_j}$-measure.
On $I_j$ $(j=4,5,7)$ the argument $s(x)$ lies in a compact subinterval of
$(1,\infty)$, where $F$ is continuous. On $\mathcal D_{10}$,
$y<K-11\alpha/5-5x/2$ gives $q(x,y)>11/5$, and $q+h/\alpha<4$; there $f$ is
continuous. Hence $k_j^0,k_j^h$ and $d_j$ are bounded and continuous off a
null set, so \eqref{eq:retB-portmanteau} applies to them.
With \eqref{eq:retB-density-limits}, it gives
\begin{equation}\label{eq:retB-direct-convergence}
 S_{j,N}^0\longrightarrow P_j,\qquad
 B_{j,N}\longrightarrow\Delta_j\qquad(j\in\mathcal H).
\end{equation}

\emph{Step 4: the ordered three-dimensional kernel.}
The cone $\mathcal C$ has affine boundary faces, so its boundary is
$m^3$-null.
The function $\Theta_*$ is piecewise affine.
Its possible jumps lie on finitely many guard hyperplanes, and its branch
switches add finitely many affine faces.
On every remaining open piece, $a_3$ and $a_3^s$ are continuous.
The denominator satisfies $t_3\ge u>0$, so the arguments of $f$ remain in a
fixed compact interval. On it $f$, extended by zero below $2$, is
continuous. Since $0\le f\le1$, $|d_G|\le K_G/u$.
Thus $d_G$ is bounded and $m^3$-almost everywhere continuous.
Apply \eqref{eq:retB-portmanteau} to the zero extension on $X_0^3$:
\begin{equation}\label{eq:retB-G-portmanteau}
 \iiint d_G\,d(m_N^3)\longrightarrow
 \iiint d_G\,d(m^3)=\Delta_G.
\end{equation}
The normalization costs at most
\begin{equation}\label{eq:retB-G-normalization-error}
 \left|B_{G,N}-\iiint d_G\,d(m_N^3)\right|
 \le\sup_{[u,\mu_*]}|r_{G,N}-1|\,\|d_G\|_\infty c_N^3
 \longrightarrow0.
\end{equation}
The cone's coordinate-order inequalities are strict, so $B_{G,N}$ counts
strictly ordered triples, as in \eqref{eq:retB-G-prime-sum}.
Hence $B_{G,N}\to\Delta_G$, and with Step~3 all limits in
\eqref{eq:retB-conversion-limit} and $\Lambda^0_{j,N}\to0$ follow.
\end{proof}

\paragraph{Whole cells, retained cells and the label cutoff.}
For the four direct rows, the arithmetic identification is
Lemma~\ref{lem:retB-lifted-direct}. For the joint row it is the proof of
\eqref{eq:retB-G-use}. Both are made at stage eight on Paper~II's actual
cells, with their inherited signed coefficients. Hence no further whole-cell
kernel is needed. For the retained part of $S_{10}$, the actual whole cells
cover $\mathcal D_{10}$ up to a strip. Lemma~\ref{lem:retB-lifted-direct}
absorbs this strip by the retained reserve \eqref{eq:retB-S10-surplus}. For
$G_1$, the retained cells are those of Paper~II, (4.19). By
Lemma~\ref{lem:retB-shifted-retention}, the shifted profile has
no positive support outside them. The arrays \eqref{eq:retB-direct-arrays}
and \eqref{eq:retB-G-array} enter only through the nonnegative errors in
$E'$, which vanish by Lemma~\ref{lem:retB-conversion}.

The array \eqref{eq:retB-G-prime-sum} carries $V_N(p_3)$ at the actual
coordinate. The arithmetic lower main term carries $V_N(E)$,
$E=N^{t_3}$, at the inherited label. Since $E\le a_3$, $V_N(E)\ge V_N(a_3)$.
In the lower part of the joint upper bound this replacement is
favourable. We nevertheless bound it absolutely and include it once in
Paper~II's peel allowance (4.57). No fixed cost is added.
By \eqref{eq:retB-Euler-limit}, uniformly for actual label/coordinate pairs in the same cell with $u/\rho\le t_3\le x_3\le1$ and $x_3<\rho t_3$,
\[
 \frac{V_N(N^{t_3})}{V_N(N^{x_3})}=\frac{x_3}{t_3}(1+o(1))\le\rho(1+o(1)).
\]
For a weight $0\le w\le1$ on the strictly ordered cone, the replacement
therefore costs at most
\begin{equation}\label{eq:retB-label-cost}
\begin{gathered}
 C_{\mathrm{cut},N}\le(\tau+o_\sigma(1))T_N,\qquad
 T_N=\frac{X_G}W\sum_{\mathcal C}
       \frac{V_N(p_3)}{(p_1-1)(p_2-1)(p_3-1)}\longrightarrow K_GI_0,\\
 I_0=\iiint_{\mathcal C}\frac{d\mathbf t}{t_1t_2t_3^2}
    =500\Bigl(1-H_{\log}+\frac{H_{\log}^2}2\Bigr)-\frac1{\mu_*},
\end{gathered}
\end{equation}
where the sum runs over the primes of \eqref{eq:retB-G-prime-sum}. The limit
follows from Steps~2--4 of the proof of Lemma~\ref{lem:retB-conversion}
applied to the bounded kernel $\mathbf 1_{\mathcal C}K_G/t_3$. We evaluate $I_0$
by \eqref{eq:retB-G-cancellation}. Since $1<H_{\log}<459/100$ and
the quadratic increases for $H_{\log}>1$,
\[
 I_0<\frac{693381}{200},\qquad
 K_GI_0\le K_{\mathrm{hi}}\frac{693381}{200}<1400.
\]
At fixed $\tau$, eventually $o_\sigma(1)\le\tau$ and $T_N\le1400$, so
$C_{\mathrm{cut},N}\le2800\tau<4000\tau$. We include $4000\tau$ once.
Paper~II's peel bound before its final rounding is
\[
 P_{\mathrm{peel}}\tau=\frac{3\cdot502002\,(459/100)^3}{1-\varepsilon}\tau
 =\frac{8090814677193}{55550}\tau
\]
\cite[(4.47), (4.56)]{PaperII}, and its printed allowance is
$P_{\mathrm{cap}}\tau=(485500000000/3333)\tau$ \cite[(4.57)]{PaperII}.
The per-layer coefficient 502002 in Paper~II, (4.53), includes the
Selberg remainder. The printed allowance (4.57) uses that coefficient;
we therefore use its slack with the Selberg remainder already included.
Exact subtraction gives
\begin{equation}\label{eq:retB-peel-room}
 P_{\mathrm{cap}}-P_{\mathrm{peel}}=\frac{2555968421}{166650},\qquad
 P_{\mathrm{cap}}-P_{\mathrm{peel}}-4000=\frac{1889368421}{166650}>0.
\end{equation}
Hence the peel and the label replacement together stay below
$P_{\mathrm{cap}}\tau<145664566.4567\tau$, the entry already contained in
$g_\tau$. The fixed deduction $F_{\mathrm{net}}$ is unchanged.

\subsection{The resulting lower comparison}

\begin{remark}[Inputs]\label{rem:retB-inputs}
The proof of Proposition~\ref{prop:retB-comparison} uses the cited
statements of Paper~I and Paper~II. They are applied to Paper~II's unchanged sources.
It uses Paper~I, Theorem~2.8, through
Fact~\ref{fact:spectral-implication}, and the prime Siegel--Walfisz theorem,
Paper~I, (2.21). It uses Lemma~\ref{hyp:ra-ms-transfer}(ii) at
Paper~I, (3.48), and
the finite enclosures of Propositions~\ref{prop:app-h6-four-profiles}
and~\ref{prop:app-m5-cap}. The optional endpoint improvement also uses
Lemma~\ref{lem:retB-six-endpoints}. The table of
Appendix~\ref{subsec:app-retB-rows} is an index, not an input.
\end{remark}

\begin{proposition}[Retained comparison]\label{prop:retB-comparison}
Fix $0<\sigma\le10^{-10}$ satisfying Paper~II, (1.27) and assume
Lemma~\ref{lem:retB-six-endpoints}.
Then \eqref{eq:retB-goal} holds, and $E'(N;\sigma)\to0$ as $N\to\infty$.
Without that lemma, \eqref{eq:retB-goal} holds with $g_6$ replaced by $0$;
the same limit for $E'(N;\sigma)$ holds.
\end{proposition}

\begin{proof}
Write each chain of $\mathcal I$, including $\mathrm{T5}$, by
\eqref{eq:retB-telescope} and each joint group by \eqref{eq:retB-joint}. Lemma~\ref{lem:retB-base-errors} covers the
differences at stages one to eight and the remaining Paper~II differences at
stages nine and ten. Their sum is at least $-(F_{\mathrm{net}}-\delta_{G1})-E_{\mathrm{van}}$.
For $j\in\mathcal H$, replace the stage-nine object by $Q^h_{j,9}=P_j+\Delta_j$.
Then Lemma~\ref{lem:retB-lifted-direct} gives $Q_{j,8}-Q^h_{j,9}\ge0$.
For $j\in\mathcal J$, Lemma~\ref{lem:retB-six-endpoints} gives the
stage-ten difference $P_j-N_j\ge0$, so the endpoint rises by $N_j-O_j$.
Only the endpoint changes, so $g_6$ enters once.
By Proposition~\ref{prop:app-h6-four-profiles},
$\sum_{j\in\mathcal H}\Delta_j>572004/10^8>5719/10^6$, which is
\eqref{eq:retB-H6-finite}.
For the joint occurrence, \eqref{eq:retB-G-use} replaces Paper~II's $G_1$
bound. Its $E^D_{G1}$ is the $G_1$ part of $F_{\mathrm{net}}+E_{\mathrm{van}}$
and is counted once.
These modifications concern disjoint profile terms, so their gains enter once.
The endpoint sum \eqref{eq:retB-base-endpoints} and the floor
\eqref{eq:retB-paperII-floor} give $D_{\mathrm{II}}+\overline J_6+2\overline H$.
Hence
\begin{equation}\label{eq:retB-sum}
\begin{aligned}
 \frac{R_{\mathrm{fix}}+S_2}{W}
 \ge{}&D_{\mathrm{II}}+g_6+\overline J_6+2\overline H
      +\frac{5719}{10^6}+M_{5,\mathrm{internal}}\\
 &-E_{\mathrm{van}}(N;\sigma)
      -\sum_{j\in\mathcal H}
          \bigl(|\Lambda^0_{j,N}|+|\Lambda_{j,N}|\bigr)-|\Lambda_{G,N}|.
\end{aligned}
\end{equation}
The nine conversion errors are nonnegative, so \eqref{eq:retB-sum} is weaker
than the same bound without them. This is \eqref{eq:retB-goal}.
By Lemmas~\ref{lem:retB-base-errors} and~\ref{lem:retB-conversion},
$E'(N;\sigma)\to0$; the threshold depends on the fixed $\sigma$.
No fixed increment is added. The only new fixed cost is the label
replacement, and by \eqref{eq:retB-peel-room} it lies inside Paper~II, (4.57).
Without Lemma~\ref{lem:retB-six-endpoints}, take $N_j=O_j$; then the
same argument gives \eqref{eq:retB-sum} with $g_6$ replaced by $0$.
\end{proof}
